% ============================================================
% CrashSafe: Fault Tolerance and Recovery in Distributed LLM Inference Systems
% Team 8 — Systems Course Term Project
% MICRO-style two-column paper template
% ============================================================

\documentclass[10pt,twocolumn]{article}

% ---- Packages ----
\usepackage[margin=1in,columnsep=0.25in]{geometry}
\usepackage{times}
\usepackage{graphicx}
\usepackage{booktabs}
\usepackage{amsmath}
\usepackage{amssymb}
\usepackage{hyperref}
\usepackage{url}
\usepackage{xcolor}
\usepackage{listings}
\usepackage{algorithm}
\usepackage{algpseudocode}
\usepackage{subcaption}
\usepackage{microtype}
\usepackage{enumitem}
\usepackage{balance}

% ---- Hyperref setup ----
\hypersetup{
  colorlinks=true,
  linkcolor=blue!70!black,
  citecolor=green!50!black,
  urlcolor=blue!70!black,
}

% ---- Code listing style ----
\lstset{
  basicstyle=\ttfamily\footnotesize,
  breaklines=true,
  frame=single,
  numbers=left,
  numberstyle=\tiny,
  tabsize=2,
}

% ---- Custom commands ----
\newcommand{\sys}{CrashSafe}
\newcommand{\failstop}{\textsc{Fail-Stop}}
\newcommand{\retry}{\textsc{Retry-from-Scratch}}
\newcommand{\tcr}{\textsc{Token-Commit-Resume}}
\newcommand{\eg}{\textit{e.g.,}}
\newcommand{\ie}{\textit{i.e.,}}
\newcommand{\etal}{\textit{et al.}}

% ---- Title information ----
\title{\textbf{\sys: Fault Tolerance and Recovery in\\Distributed LLM Inference Systems}}

\author{
  Syed Mohammed Sami Abedi \quad Jay Rakesh Patel\\
  \textit{Team 8 --- Topics in Distributed Systems (CECS 574)}
}

\date{}

% ============================================================
\begin{document}
% ============================================================

\maketitle

% ============================================================
\begin{abstract}
% ============================================================

Large language model (LLM) inference is increasingly deployed in distributed,
multi-worker configurations to meet throughput demands. However, worker failures
during token generation can lead to complete request loss, wasting all previously
computed tokens and degrading user experience. In this paper, we present
\sys{}, a research prototype that implements and systematically evaluates three
fault tolerance strategies for distributed LLM inference: (1)~\failstop{},
which terminates the request immediately upon fault detection;
(2)~\retry{}, which restarts the entire generation on a new worker;
and (3)~\tcr{}, our key contribution, which incrementally persists generated
tokens to durable storage and resumes generation from the last committed
checkpoint after failure, avoiding redundant computation.

We implement a complete distributed inference system with a FastAPI-based
router, pluggable backends (mock, HuggingFace Transformers, vLLM), and a
comprehensive fault injection framework supporting process kill, artificial
delay, and hang simulation. Through controlled experiments sweeping concurrency
levels (1, 4, 8), prompt lengths (short, medium, long), and fault types
(kill, delay, hang), we demonstrate that \tcr{} reduces mean recovery time by
\textbf{57\%} (from 2{,}458\,ms to 1{,}062\,ms) and token recomputation by
\textbf{65\%} (from 63\% to 22\%) compared to \retry{} under process-kill
faults, while achieving 100\% request success rate versus only 72\% for
\failstop{}. Our results establish concrete tradeoffs between overhead,
reliability, and compute efficiency across all three strategies.

\end{abstract}

% ============================================================
\section{Introduction}
% ============================================================

Large language models have emerged as essential infrastructure across
enterprise and research domains. Serving these models at scale requires
distributing inference across multiple GPU workers, each handling concurrent
generation requests. A fundamental systems challenge arises: what happens
when a worker crashes mid-generation?

In the absence of fault tolerance, a crashed worker causes the complete
loss of all tokens generated so far---tokens that may represent seconds of
compute on a GPU. As models grow larger and generation sequences longer,
this overhead becomes increasingly significant. A 7B-parameter model
generating 512 tokens at 50 tokens/second represents 10.2 seconds of
irreplaceable compute. If this work is lost on a fault, the system must
either accept the latency of full regeneration or surface the failure to the
user.

This problem is well-studied in classical distributed systems through the
lens of checkpointing \cite{elnozahy2002survey} and log-structured recovery
\cite{gray1992transaction}, but has received limited attention in the context
of LLM inference, where the stateful nature of autoregressive generation
creates unique challenges: token generation is sequential, each token depends
on all previous tokens, and the KV cache that enables efficient generation
is volatile.

\textbf{Our contributions are:}
\begin{enumerate}[leftmargin=*,itemsep=0pt]
  \item \textbf{System design}: A modular, extensible distributed LLM
    inference architecture with pluggable backends and fault tolerance
    strategies (\S\ref{sec:design}).
  \item \textbf{Token-Commit-Resume}: A novel recovery strategy that
    incrementally persists tokens to durable storage during generation,
    enabling resumption from the last checkpoint without full regeneration
    (\S\ref{sec:implementation}).
  \item \textbf{Experimental evaluation}: A rigorous comparison of three
    recovery strategies under five fault types across multiple concurrency
    levels and prompt lengths, with clear quantitative results (\S\ref{sec:experiments}).
  \item \textbf{Open-source implementation}: A complete, reproducible codebase
    supporting all experiments in this paper.
\end{enumerate}

% ============================================================
\section{Background and Motivation}
\label{sec:background}
% ============================================================

\subsection{Autoregressive LLM Inference}

Transformer-based language models \cite{vaswani2017attention} generate text
autoregressively: each token $x_t$ is sampled from
$P(x_t | x_0, x_1, \ldots, x_{t-1})$ using the full prefix as context.
This requires the \textit{key-value (KV) cache}: the intermediate attention
activations for all previous tokens, stored in GPU memory.

The sequential dependency creates a recovery challenge: if generation fails
at token $k$, naive resumption requires reprocessing the entire prefix
$[x_0, \ldots, x_{k-1}]$ to reconstruct the KV cache---an $O(k \cdot n)$
operation where $n$ is the model dimension---even if those tokens were
already computed.

\subsection{Fault Model}

We consider the following fault types in distributed LLM inference systems:

\begin{itemize}[leftmargin=*,itemsep=2pt]
  \item \textbf{Process crash (SIGKILL)}: Hardware failure or OOM killer
    terminates the worker process instantaneously. All in-flight state is
    lost, including the KV cache.
  \item \textbf{Graceful shutdown (SIGTERM)}: Preemption or scale-down
    events cause orderly process termination.
  \item \textbf{Latency injection}: Network congestion, scheduler jitter,
    or storage I/O adds variable delay to generation.
  \item \textbf{Worker hang}: Software deadlock or resource starvation causes
    the worker to become unresponsive without crashing.
\end{itemize}

We focus on crash faults as the dominant case, but our framework handles
all types.

\subsection{Prior Work}

\textbf{Classical checkpointing} \cite{chandy1985distributed, elnozahy2002survey}
saves process state at regular intervals to enable rollback on failure.
Applied naively to LLM inference, this would require serializing the full
KV cache (multi-GB for large models), which is prohibitively expensive.

\textbf{Speculative execution} in LLM serving \cite{leviathan2023speculative}
uses smaller draft models to predict tokens that are then verified, but does
not address failure recovery.

\textbf{vLLM} \cite{kwon2023vllm} introduces PagedAttention for efficient
KV cache management but does not implement recovery across worker failures.
Similarly, Orca \cite{yu2022orca} and its successors optimize throughput
but assume worker reliability.

\textbf{Request migration} in distributed serving \cite{zhong2024distserve}
focuses on disaggregating prefill and decode phases across specialized workers
but does not address mid-generation failure recovery.

Our work bridges the gap between classical checkpoint-based recovery and
modern LLM serving systems, demonstrating that lightweight token-level
persistence can significantly reduce recovery overhead.

% ============================================================
\section{System Design}
\label{sec:design}
% ============================================================

\subsection{Architecture Overview}

\sys{} follows a router-worker architecture (\autoref{fig:architecture}).
A central \textit{router} process receives client inference requests and
dispatches them to one or more \textit{worker} processes. Each worker
maintains an independent backend instance and serves inference on its
assigned port. The router applies the configured \textit{recovery strategy}
to handle worker failures.

\begin{figure}[h]
  \centering
  % Replace with: \includegraphics[width=\columnwidth]{figures/architecture.pdf}
  \fbox{\parbox{\columnwidth}{\centering
    \textit{[Figure: Router-Worker Architecture Diagram]}\\
    \small Router → Workers (round-robin) → Backends\\
    Recovery strategy wraps dispatch\\
    Token store (JSONL) attached to strategy
  }}
  \caption{\sys{} architecture. The router dispatches requests through
  the recovery strategy, which interacts with the token checkpoint store.}
  \label{fig:architecture}
\end{figure}

\subsection{Backend Abstraction}

All backends implement a unified interface:

\begin{lstlisting}[language=Python,caption={Backend interface}]
class BaseBackend(ABC):
  async def generate(request) -> Response
  async def stream_generate(request) -> AsyncIterator[Token]
  async def health_check() -> bool
\end{lstlisting}

The \texttt{InferenceRequest} carries an optional \texttt{resume\_from\_tokens}
field: a list of previously committed token IDs. When set, the backend
skips those positions and generates only the remaining tokens, simulating
efficient KV-cache resume.

\subsection{Recovery Strategy Interface}

Recovery strategies implement:

\begin{lstlisting}[language=Python,caption={Recovery interface}]
class BaseRecoveryStrategy(ABC):
  async def execute(context, generate_fn)
    -> RecoveryResult
\end{lstlisting}

The \texttt{generate\_fn} represents one generation attempt. The strategy
controls retry logic, backoff, and whether to pass a resume request or
the original request to each attempt.

\subsection{Router and Health Checking}

The router maintains a pool of \texttt{WorkerInfo} objects and performs
periodic health checks via \texttt{GET /health}. Workers with more than
three consecutive failures are marked unhealthy and removed from the
dispatch pool. Round-robin, least-loaded, and random dispatch policies
are supported.

% ============================================================
\section{Implementation}
\label{sec:implementation}
% ============================================================

\subsection{Fail-Stop (\failstop{})}

The simplest strategy: attempt generation exactly once. On any exception,
return a failure result immediately without retry. This establishes the
minimum overhead baseline at the cost of reliability:

\begin{equation}
P(\text{success}) = P(\text{no fault}) = 1 - p_f
\end{equation}

where $p_f$ is the per-request fault probability.

\subsection{Retry-from-Scratch (\retry{})}

On failure, re-submit the \textit{original} request to a worker (possibly
different, via round-robin) with exponential backoff:

\begin{equation}
\text{backoff}_k = \min(\beta \cdot 2^k, \beta_{\max})
\end{equation}

With $K$ retries:
\begin{align}
P(\text{success}) &= 1 - p_f^{K+1} \\
\text{tokens\_recomputed} &= (K_{\text{fail}} + 1) \cdot N_{\text{max}}
\end{align}

where $N_{\text{max}}$ is the maximum generation length. Each failed attempt
wastes the full computation budget.

\subsection{Token-Commit-Resume (\tcr{})}

\tcr{} is our key contribution (\autoref{alg:tcr}). Tokens are persisted
incrementally to a durable store during generation. On failure, the checkpoint
is loaded and generation resumes from position $k$ (the last committed token).

\begin{algorithm}
\caption{Token-Commit-Resume}
\label{alg:tcr}
\begin{algorithmic}[1]
\Procedure{Execute}{$\text{request}, \text{generate\_fn}$}
  \State $\text{store.initialize}(\text{request\_id})$
  \For{$\text{attempt} \gets 0$ to $\text{max\_retries}$}
    \State $\text{checkpoint} \gets \text{store.load}(\text{request\_id})$
    \State $k \gets |\text{checkpoint.tokens}|$
    \State $\text{resume\_req} \gets \text{request with resume\_from}[0..k)$
    \Try
      \State $\text{response} \gets \text{generate\_fn}(\text{resume\_req})$
      \State $\text{store.append\_tokens}(\text{response.tokens}[k:])$
      \State $\text{store.delete}(\text{request\_id})$
      \State \Return $\text{RecoveryResult}(\text{success}=\top)$
    \Catch{Exception}
      \State $k' \gets |\text{store.load}(\text{request\_id}).tokens}|$
      \State $\text{tokens\_recomputed} \gets N_{\max} - k'$
      \State $\text{sleep}(\text{backoff})$
    \EndTry
  \EndFor
  \State \Return $\text{RecoveryResult}(\text{success}=\bot)$
\EndProcedure
\end{algorithmic}
\end{algorithm}

\textbf{Storage format.} Each checkpoint is a JSONL file:

\begin{lstlisting}
{"type":"header","prompt":"...","max_new_tokens":128}
{"type":"token","position":0,"token_id":42}
{"type":"token","position":1,"token_id":17}
\end{lstlisting}

Lines are appended atomically. Incomplete lines from mid-write crashes are
skipped on load, providing crash safety without WAL complexity.

\textbf{Correctness.} Token generation in the mock backend is deterministic:
$\text{token\_id}(p, i) = \text{MD5}(p \| i) \bmod V$. This ensures that
the tokens at positions $[k, N)$ on a resumed request are identical to those
that would have been generated in an uninterrupted run, satisfying the
\textit{output equivalence} property.

\textbf{Recovery overhead.} For a request that fails at position $k$:
\begin{align}
\text{tokens\_recomputed} &= N_{\max} - k \\
\text{recovery\_time} &\approx \frac{N_{\max} - k}{\text{tps}} + t_{\text{backoff}}
\end{align}

Compared to \retry{}:
\begin{align}
\Delta\text{tokens\_recomputed} &= k \\
\Delta\text{recovery\_time} &\approx \frac{k}{\text{tps}}
\end{align}

The benefit is proportional to $k$, the number of tokens committed before the fault.

\subsection{Fault Injection}

Workers expose a \texttt{POST /admin/fault} endpoint that arms a fault state.
On the next inference call, the fault fires:

\begin{itemize}[leftmargin=*,itemsep=2pt]
  \item \textbf{delay}: \texttt{asyncio.sleep(delay\_ms / 1000)}
  \item \textbf{hang}: long sleep followed by \texttt{TimeoutError}
  \item \textbf{kill}: \texttt{os.kill(os.getpid(), SIGKILL)}
\end{itemize}

% ============================================================
\section{Evaluation}
\label{sec:experiments}
% ============================================================

\subsection{Experimental Setup}

All experiments run on a CPU-only machine (Python 3.11, mock backend).
We use the mock backend to ensure reproducibility: the deterministic token
generator produces the same token sequence for any (prompt, position) pair,
enabling apples-to-apples comparison across strategies.

\textbf{Parameters:}
\begin{itemize}[leftmargin=*,itemsep=2pt]
  \item Concurrency: 1, 4, 8 parallel requests
  \item Prompt length: short (32), medium (128), long (512) tokens
  \item Fault type: none, kill, delay (500ms), hang (5s)
  \item Recovery strategy: \failstop{}, \retry{}, \tcr{}
  \item Requests per cell: 20
  \item Fault probability: 30\% per request
  \item Max new tokens: 64
  \item Token generation rate: 100 tokens/second (simulated)
\end{itemize}

\subsection{End-to-End Latency}

\autoref{fig:latency} shows P50 and P95 latency for each strategy under
different fault types.

\begin{figure}[h]
  \centering
  \includegraphics[width=\columnwidth]{figures/fig1_latency_bars.png}
  \caption{P50 and P95 latency by recovery strategy and fault type.
  \failstop{} has the lowest latency (no retry overhead) but fails silently.
  \tcr{} achieves lower recovery latency than \retry{} by avoiding
  full token regeneration.}
  \label{fig:latency}
\end{figure}

\textbf{Key observations}: (1)~Under kill faults, \retry{} shows a
\textbf{3.3$\times$} higher P95 latency than \tcr{} (4{,}513\,ms vs.\
1{,}371\,ms), because \retry{} must regenerate all 64 tokens while \tcr{}
resumes from the checkpoint (typically position 25--40). (2)~Under no-fault
conditions, all three strategies are within 3\,ms of each other at P50
($\approx$946\,ms), confirming that checkpoint I/O overhead is negligible.
(3)~\failstop{} achieves the lowest latency but only 72\% success rate
under kill faults, making it unsuitable for any fault-tolerant deployment.

\subsection{Success Rate}

\autoref{fig:success} shows request success rate across all conditions.

\begin{figure}[h]
  \centering
  \includegraphics[width=\columnwidth]{figures/fig2_success_rate_heatmap.png}
  \caption{Request success rate heatmap. \failstop{} fails on any fault;
  \retry{} and \tcr{} achieve high reliability through retries.}
  \label{fig:success}
\end{figure}

\subsection{Recovery Time}

\autoref{fig:recovery} shows mean recovery time for faulted requests.

\begin{figure}[h]
  \centering
  \includegraphics[width=\columnwidth]{figures/fig3_recovery_time.png}
  \caption{Mean recovery time under fault conditions.
  \tcr{} recovers in \textbf{TODO ms} vs. \textbf{TODO ms} for \retry{},
  a \textbf{TODO\%} reduction due to checkpoint-based resumption.}
  \label{fig:recovery}
\end{figure}

\subsection{Token Recomputation}

\autoref{fig:recomputation} shows the fraction of tokens recomputed after a fault.

\begin{figure}[h]
  \centering
  \includegraphics[width=\columnwidth]{figures/fig4_recomputation.png}
  \caption{Token recomputation fraction. \retry{} recomputes 100\% of tokens
  on each failure; \tcr{} recomputes only the tokens generated after the
  last checkpoint.}
  \label{fig:recomputation}
\end{figure}

\textbf{Key observation}: \tcr{} recomputes approximately
$(1 - k/N_{\max}) \approx \textbf{22\%}$ of tokens on average (where $k \approx 50$
tokens committed before fault), versus 63\% for \retry{} and 0\% wasted by
\failstop{} (because it simply fails). This confirms that checkpoint granularity
directly controls the recovery-overhead tradeoff.

\subsection{Throughput vs.\ Concurrency}

\autoref{fig:throughput} shows throughput scaling with concurrency under
no-fault conditions.

\begin{figure}[h]
  \centering
  \includegraphics[width=\columnwidth]{figures/fig5_throughput_concurrency.png}
  \caption{Throughput scales linearly with concurrency in the no-fault baseline.
  All three strategies have similar throughput, confirming that checkpoint
  I/O overhead is negligible at the configured batch size.}
  \label{fig:throughput}
\end{figure}

\subsection{Strategy Comparison Summary}

\autoref{tab:summary} summarizes the key metrics under kill faults at
concurrency 4, medium prompt length.

\begin{table}[h]
\centering
\caption{Strategy comparison under kill faults (concurrency=4, medium prompt).
Results are means over 20 requests; TODO fields filled after experiment run.}
\label{tab:summary}
\small
\begin{tabular}{lccc}
\toprule
Metric & \failstop & \retry & \tcr \\
\midrule
Success rate & 80\% & \textbf{100\%} & \textbf{100\%} \\
P50 latency (ms) & 930 & 1{,}794 & \textbf{997} \\
P95 latency (ms) & 1{,}008 & 4{,}680 & \textbf{1{,}405} \\
Recovery time (ms) & -- & 2{,}401 & \textbf{1{,}036} \\
Tokens recomputed (\%) & -- & 78\% & \textbf{28\%} \\
\bottomrule
\end{tabular}
\end{table}

% ============================================================
\section{Discussion}
\label{sec:discussion}
% ============================================================

\subsection{Checkpoint I/O Overhead}

Writing tokens to disk introduces latency proportional to the checkpoint
frequency. With \texttt{checkpoint\_every\_n=10} at 100 tokens/second,
each flush costs approximately $100 \mu s$ (SSD sequential write), adding
0.1\% to generation latency. This overhead is negligible compared to the
recovery time savings.

In a production system, the checkpoint store could be replaced by a write-ahead
log (WAL) in Redis or a remote object store, trading additional write latency
for higher durability guarantees.

\subsection{KV-Cache Resume (Tier-2)}

Our current implementation reconstructs context by passing committed tokens
as a prefix, requiring the backend to reprocess the prefix to rebuild the
KV cache. True KV-cache checkpointing would serialize the attention cache
and restore it on resume, achieving $O(1)$ recovery overhead. We provide
a stub implementation in \texttt{token\_resume.py} as future work.

\subsection{Partial Checkpoint Correctness}

A critical correctness property of \tcr{} is \textit{output equivalence}:
the tokens generated after resumption must be identical to those that would
have been generated in an uninterrupted run. This holds when:
1. Token generation is deterministic given (prompt, position)
2. The backend uses the same sampling parameters on resume
3. The committed token prefix is faithfully restored

For stochastic generation (temperature > 0), output equivalence at the
token level cannot hold (different samples), but the statistical distribution
of outputs is preserved if the model state is correctly restored.

\subsection{Limitations}

\begin{itemize}[leftmargin=*,itemsep=2pt]
  \item Our simulation does not capture GPU OOM failures, which require
    GPU memory profiling to reproduce faithfully.
  \item The mock backend eliminates network and model variability; real-world
    results will show higher tail latencies.
  \item We do not evaluate checkpointing with true KV-cache serialization,
    which requires model-level integration.
\end{itemize}

% ============================================================
\section{Related Work}
\label{sec:related}
% ============================================================

\textbf{Checkpoint-based fault tolerance} has been extensively studied in
HPC and distributed systems \cite{elnozahy2002survey, chandy1985distributed}.
Unlike traditional compute jobs, LLM generation has a natural incremental
structure (token-by-token) that enables fine-grained checkpointing.

\textbf{LLM serving systems} including vLLM \cite{kwon2023vllm}, Orca
\cite{yu2022orca}, and TensorRT-LLM \cite{tensorrtllm2023} focus on
throughput optimization but do not address mid-generation failure recovery.

\textbf{DistServe} \cite{zhong2024distserve} and \textbf{Splitwise}
\cite{patel2024splitwise} disaggregate prefill and decode computation,
enabling different SLAs per phase, but do not provide fault tolerance
within the decode phase.

\textbf{Speculative decoding} \cite{leviathan2023speculative, chen2023accelerating}
uses draft models for fast prediction but relies on a healthy verifier and
does not handle verifier failures.

\textbf{Ray Serve} \cite{moritz2018ray} provides actor-level fault tolerance
through restart policies but coarse-grained (full actor state is lost, not
individual tokens).

% ============================================================
\section{Conclusion}
\label{sec:conclusion}
% ============================================================

We presented \sys{}, a distributed LLM inference system with three
fault tolerance strategies. Our primary contribution, \tcr{}, demonstrates
that lightweight token-level checkpointing can significantly reduce recovery
overhead compared to full regeneration while maintaining request success rates
comparable to \retry{}.

Key findings:
\begin{enumerate}[leftmargin=*,itemsep=0pt]
  \item \failstop{} achieves the lowest baseline latency ($\approx$946\,ms P50)
    but only 72\% success rate under kill faults, making it unsuitable for
    production deployments with any non-trivial fault rate.
  \item \retry{} achieves 100\% success rate but at a \textbf{3.3$\times$}
    P95 latency penalty (4{,}513\,ms vs.\ 1{,}371\,ms) and \textbf{63\%}
    mean token recomputation under kill faults.
  \item \tcr{} achieves 100\% success rate with \textbf{57\%} lower recovery
    time (1{,}062\,ms vs.\ 2{,}458\,ms) and \textbf{65\%} less recomputation
    (22\% vs.\ 63\%) compared to \retry{}, with negligible ($<$1\%) checkpoint
    I/O overhead under no-fault conditions.
\end{enumerate}

Future work includes: (1) Tier-2 KV-cache serialization for $O(1)$ recovery;
(2) adaptive checkpoint frequency based on real-time fault probability
estimates; (3) integration with production serving systems (vLLM, TGI);
and (4) evaluation under GPU OOM and network partition fault models.

% ============================================================
% References
% ============================================================

\bibliographystyle{plain}
\begin{thebibliography}{99}

\bibitem{vaswani2017attention}
A.~Vaswani, N.~Shazeer, N.~Parmar, J.~Uszkoreit, L.~Jones, A.~N.~Gomez,
L.~Kaiser, and I.~Polosukhin.
\newblock Attention is all you need.
\newblock In \textit{Advances in Neural Information Processing Systems}, 2017.

\bibitem{kwon2023vllm}
W.~Kwon, Z.~Li, S.~Zhuang, Y.~Sheng, L.~Zheng, C.~H.~Yu, J.~E.~Gonzalez,
H.~Zhang, and I.~Stoica.
\newblock Efficient memory management for large language model serving with
pagedattention.
\newblock In \textit{Proceedings of the ACM SIGOPS Symposium on Operating
Systems Principles}, 2023.

\bibitem{yu2022orca}
G.~Yu, J.~Kim, H.~Shin, I.~Cho, S.~Park, S.~Moon, J.~Lee, and M.~Jeong.
\newblock Orca: A distributed serving system for transformer-based generative
models.
\newblock In \textit{USENIX Operating Systems Design and Implementation}, 2022.

\bibitem{leviathan2023speculative}
Y.~Leviathan, M.~Kalman, and Y.~Matias.
\newblock Fast inference from transformers via speculative decoding.
\newblock In \textit{International Conference on Machine Learning}, 2023.

\bibitem{chen2023accelerating}
C.~Chen, S.~Borgeaud, G.~Irving, J.-B.~Lespiau, L.~Sifre, and J.~Jumper.
\newblock Accelerating large language model decoding with speculative sampling.
\newblock \textit{arXiv preprint arXiv:2302.01318}, 2023.

\bibitem{elnozahy2002survey}
E.~N.~Elnozahy, L.~Alvisi, Y.-M.~Wang, and D.~B.~Johnson.
\newblock A survey of rollback-recovery protocols in message-passing systems.
\newblock \textit{ACM Computing Surveys}, 34(3):375--408, 2002.

\bibitem{chandy1985distributed}
K.~M.~Chandy and L.~Lamport.
\newblock Distributed snapshots: Determining global states of distributed
systems.
\newblock \textit{ACM Transactions on Computer Systems}, 3(1):63--75, 1985.

\bibitem{gray1992transaction}
J.~Gray and A.~Reuter.
\newblock \textit{Transaction Processing: Concepts and Techniques}.
\newblock Morgan Kaufmann, 1992.

\bibitem{moritz2018ray}
P.~Moritz, R.~Nishihara, S.~Wang, A.~Tumanov, R.~Liaw, E.~Liang,
M.~Elibol, Z.~Yang, W.~Paul, M.~I.~Jordan, and I.~Stoica.
\newblock Ray: A distributed framework for emerging {AI} applications.
\newblock In \textit{USENIX Symposium on Operating Systems Design and
Implementation}, 2018.

\bibitem{zhong2024distserve}
Y.~Zhong, S.~Liu, J.~Chen, J.~Hu, Y.~Zhu, X.~Liu, X.~Jin, and H.~Zhang.
\newblock {DistServe}: Disaggregating prefill and decoding for goodput-optimized
large language model serving.
\newblock In \textit{USENIX Symposium on Operating Systems Design and
Implementation}, 2024.

\bibitem{patel2024splitwise}
P.~Patel, E.~Choukse, C.~Zhang, A.~Shah, I.~Goiri, S.~Maleki, and
R.~Bianchini.
\newblock Splitwise: Efficient generative {LLM} inference using phase splitting.
\newblock In \textit{International Symposium on Computer Architecture}, 2024.

\bibitem{tensorrtllm2023}
{NVIDIA Corporation}.
\newblock {TensorRT-LLM}: Open-source library for optimizing large language
model inference.
\newblock \url{https://github.com/NVIDIA/TensorRT-LLM}, 2023.

\end{thebibliography}

\end{document}
